\chapter{質問紙の全項目}
\label{app:questionnaire}
OshiFitで用いる質問紙（全67項目）を示す．Q42は注意確認項目である．

\section*{あなたご自身と推し活について}
\noindent 以下の質問には「主たる推し」（いま最も時間を使っている推しを1人・1組）について答えてください。\\
\begin{longtable}{@{}p{0.08\linewidth}p{0.86\linewidth}@{}}
\toprule 番号 & 項目 \\ \midrule \endhead \bottomrule \endfoot
Q1 & 年齢 \\
Q2 & 性別\newline {\footnotesize（男性，女性，その他，回答しない）} \\
Q3 & 職業\newline {\footnotesize（学生，会社員・公務員，自営業・フリーランス，パート・アルバイト，主婦・主夫，その他）} \\
Q4 & 主たる推しのジャンル\newline {\footnotesize（実在のアイドル，アニメ・漫画・ゲームのキャラクター，VTuber・配信者，俳優・舞台役者，声優，バンド・アーティスト，スポーツ選手，その他）} \\
Q5 & その推しを応援し始めてからの期間\newline {\footnotesize（半年未満，半年〜1年未満，1年〜3年未満，3年〜5年未満，5年以上）} \\
Q6 & 推し活に使う1か月あたりの金額（目安）\newline {\footnotesize（5,000円未満，5,000円〜1万円未満，1万円〜3万円未満，3万円〜5万円未満，5万円以上，答えたくない）} \\
\end{longtable}
\section*{あなた自身の性格について（TIPI-J）}
\noindent 以下の項目はどの程度当てはまりますか。直感でお答えください。\\
\noindent {\small 回答：1：全く違うと思う／2：少し違うと思う／3：どちらでもない／4：少しそう思う／5：強くそう思う}
\begin{longtable}{@{}p{0.08\linewidth}p{0.86\linewidth}@{}}
\toprule 番号 & 項目 \\ \midrule \endhead \bottomrule \endfoot
Q7 & 活発で外向的だと思う \\
Q8 & 他人に不満をもち、もめごとを起こしやすいと思う \\
Q9 & しっかりしていて、自分に厳しいと思う \\
Q10 & 心配性で、うろたえやすいと思う \\
Q11 & 新しいことが好きで、変わった考えをもつと思う \\
Q12 & 寡黙で静かだと思う \\
Q13 & 思いやりがあり、温かいと思う \\
Q14 & だらしなく、うっかりしていると思う \\
Q15 & 冷静で、気分が安定していると思う \\
Q16 & 発想力が乏しく、平凡だと思う \\
\end{longtable}
\section*{推し活の中身（割合）}
\noindent 回数ではなく、あなたの推し活全体の中で、それぞれの活動がどのくらいの割合を占めているかをお答えください。\\
\noindent {\small 回答：1：まったくしない／2：ごく一部（たまにする程度）／3：一部（いくつかある活動の一つ）／4：大きな部分（主な活動の一つ）／5：中心（推し活のほとんどがこれ）}
\begin{longtable}{@{}p{0.08\linewidth}p{0.86\linewidth}@{}}
\toprule 番号 & 項目 \\ \midrule \endhead \bottomrule \endfoot
Q17 & 推しに関係する音楽を聴いている \\
Q18 & 推しに関連する動画（YouTube、MVなど）を視聴している \\
Q19 & 推しのライブやイベントに現地参加している \\
Q20 & 推しのCDやBlu-rayなどを購入している \\
Q21 & 推しのグッズ（アクスタ、タオル等）を購入している \\
Q22 & 推しが出演するテレビ番組やラジオをチェックしている \\
Q23 & 推しに関連する雑誌や書籍を購入している \\
Q24 & 推しのSNS（X、Instagram等）をフォロー・チェックしている \\
Q25 & SNSで推しに関する情報を発信・共有している \\
Q26 & 推しの聖地巡礼（ロケ地めぐり等）を行っている \\
Q27 & ファンクラブや有料会員サービスに入会している \\
Q28 & 推しのコラボカフェやポップアップストアに行っている \\
\end{longtable}
\section*{推し活を通じた気持ちと生活の変化}
\noindent 推し活をしている中でのあなたの気持ちや生活について、どの程度当てはまりますか。\\
\noindent {\small 回答：1：全く当てはまらない／2：ほとんど当てはまらない／3：あまり当てはまらない／4：どちらともいえない／5：やや当てはまる／6：かなり当てはまる／7：非常に当てはまる}
\begin{longtable}{@{}p{0.08\linewidth}p{0.86\linewidth}@{}}
\toprule 番号 & 項目 \\ \midrule \endhead \bottomrule \endfoot
Q29 & 推し活をしていると、日々のストレスが解消されると感じる \\
Q30 & 推しの存在が、毎日の生きがい（活力）になっている \\
Q31 & 辛いことがあっても、推しのことを考えると乗り越えられる \\
Q32 & 推し活を通じて、自分に自信が持てるようになった \\
Q33 & 推し活をしている時間は、時間を忘れて没頭できる \\
Q34 & 推し活のコミュニティや活動の中で、自分の居場所があると感じる \\
Q35 & 推しのために、規則正しい生活（早起きなど）を心がけている \\
Q36 & 推し活の時間を確保するために、スケジュール管理が上手くなった \\
Q37 & 推しに見合うファンでいるために、健康や体型維持に気を使っている \\
Q38 & 推し活資金を確保するために、無駄遣いを減らし節約している \\
Q39 & 日々の生活の中で「幸せ」を感じる瞬間が増えた \\
Q40 & 情報収集能力が向上した \\
Q41 & 推し活を通じて、新しい友人（推し友）ができた \\
Q42 & この質問には必ず「2　ほとんど当てはまらない」を選んでください \\
Q43 & 家族やリアルな友人との間で、推し活の話を通じてコミュニケーションが円滑になった \\
Q44 & 推し活のために睡眠時間を削り、日常生活に支障が出ることがある \\
Q45 & 推し活にお金を使いすぎて、生活費や将来の貯蓄に不安を感じることがある \\
Q46 & 予算を超えてグッズを買ったり、ガチャを回したりして後悔することがよくある \\
Q47 & イベント参加や情報追跡を「しなければならない」という義務感（プレッシャー）を感じている \\
Q48 & SNSで他のファン（同担など）の活動を見て、嫉妬したり落ち込んだりすることがある \\
Q49 & ファンのコミュニティ内の人間関係に疲れを感じている \\
Q50 & 推し活が原因で、家族やリアルな友人との関係が悪くなったことがある \\
Q51 & 推し活をしていないと不安になったり、気持ちが不安定になったりすることがある \\
Q52 & 推し活のために、ファッションや身だしなみ（メイク、美容院など）への意識が高まった \\
Q53 & 推しを表現するために、創作スキル（イラスト、動画、文章など）が向上した \\
Q54 & 誰かと協力して応援する（企画参加など）楽しさを知った \\
\end{longtable}
\section*{推し活をしたくなるとき}
\noindent 「推し活をしたい・推しにお金や時間を使いたい」と強く感じるのは、どのようなときですか。\\
\noindent {\small 回答：1：全く当てはまらない／2：ほとんど当てはまらない／3：あまり当てはまらない／4：どちらともいえない／5：やや当てはまる／6：かなり当てはまる／7：非常に当てはまる}
\begin{longtable}{@{}p{0.08\linewidth}p{0.86\linewidth}@{}}
\toprule 番号 & 項目 \\ \midrule \endhead \bottomrule \endfoot
Q55 & 仕事や学校などで嫌なことがあり、ストレスから逃れたいとき \\
Q56 & 日常生活で孤独感や不安を感じたとき \\
Q57 & 何かを頑張った後の、自分へのご褒美として \\
Q58 & 疲れているとき、癒やしや元気をもらいたいとき \\
Q59 & SNS等で、他のファンが盛り上がっている（買っている）のを見たとき \\
Q60 & 気分に関わらず、毎日の生活の一部（習慣）として \\
Q61 & テンションを上げたいとき \\
\end{longtable}
\section*{最近2週間の気分（WHO-5）}
\noindent 最近2週間のあなたの状態に最も近いものを選んでください。\\
\noindent {\small 回答：0：まったくない／1：ほんのたまに／2：半分以下の期間を／3：半分以上の期間を／4：ほとんどいつも／5：いつも}
\begin{longtable}{@{}p{0.08\linewidth}p{0.86\linewidth}@{}}
\toprule 番号 & 項目 \\ \midrule \endhead \bottomrule \endfoot
W1 & 明るく、楽しい気分で過ごした \\
W2 & 落ち着いた、リラックスした気分で過ごした \\
W3 & 意欲的で、活動的に過ごした \\
W4 & ぐっすりと休め、気持ちよく目覚めた \\
W5 & 日常生活の中に、興味のあることがたくさんあった \\
\end{longtable}
\section*{最後に}
\begin{longtable}{@{}p{0.08\linewidth}p{0.86\linewidth}@{}}
\toprule 番号 & 項目 \\ \midrule \endhead \bottomrule \endfoot
Q62 & 推し活に関する悩みや、ストレスがたまる・幸せを感じる具体的な場面があれば自由にお書きください（任意） \\
\end{longtable}

\chapter{提案プランの詳細}
\label{app:plans}
表\ref{tab:plans}の11種類のプランについて，具体的な実行方法を示す．
\begin{longtable}{@{}p{0.3\linewidth}p{0.64\linewidth}@{}}
\toprule プラン & 実行方法 \\ \midrule \endhead \bottomrule \endfoot
SNSは23時まで．そのあとはお休みモード & 通知を23時にオフにし，翌朝まとめて見る \\
リアタイは週2回まで，ほかはアーカイブで & 見逃しは翌日の決めた時間に見る \\
今月のグッズ予算を先に決める & 予算を決め，超えたら翌月に回す \\
欲しいグッズは1日おいてから決める & カートに入れて24時間後にもう一度見る \\
今週追うものを3つに絞る & 全部追わなくていいと決めてリストにする \\
「行かない回」を1つ決めておく & 遠征やイベントは無理のない回だけ選ぶ \\
比べてしまう言葉をミュートする & 「当選」「〇〇円分」などをミュート語に登録する \\
SNSを見るのは1日2回・各15分まで & 見る時間を決め，それ以外は通知をオフにする \\
友達と同時視聴の日をつくる & 週末に1回，感想を話しながら見る \\
感想メモや創作で推しを残す & 視聴後に3行の感想，または絵・文章を残す \\
寝る前15分は推しの曲だけの時間 & 画面を見ずに音楽だけで過ごす \\
\end{longtable}

\chapter{週ごとの振り返りの項目}
\label{app:weekly}
\todo{要記入：最終版の文言に差し替える}
\begin{enumerate}
\item 今週，選んだプランを実行しましたか（行った／一部行った／行わなかった）
\item 推し疲れ4項目（睡眠への支障，お金の不安，義務感，SNSでの比較）（7件法）
\item 今週の推し活中の幸福感（7件法）
\item 先週と比べたストレスの変化（減った／変わらない／増えた）
\item 先週と比べた幸福感の変化（上がった／変わらない／下がった）
\end{enumerate}

\chapter{終了時調査の項目}
\label{app:final}
\todo{要記入：WHO-5を含む終了時調査20項目と，群B・C向けの気づき・使いやすさ・有用性7項目の最終版を記載する}

\chapter{分析の手順}
\label{app:analysis}
第3章の分析と第5章の分析は，Pythonのスクリプトとして実装した．表\ref{tab:analysis}に，分析の手順を示す．

\begin{table}[h]
\centering
\jetcaption{分析の手順}{Analysis steps}
\label{tab:analysis}
\small
\begin{tabular}{@{}p{0.05\linewidth}p{0.22\linewidth}p{0.33\linewidth}p{0.3\linewidth}@{}}
\toprule
手順 & 内容 & 手法 & 確かめること \\
\midrule
0 & データの品質確認 & 注意確認項目，同一回答の除外 & 分析に使える回答だけを残す \\
1 & 得点化と信頼性 & 逆転項目の処理，Cronbachの$\alpha$ & 尺度がまとまっているか \\
2 & 測定モデル & 確認的因子分析（充実と負担の2因子） & 楽しさと疲れが別の因子として測れているか \\
3 & 類型化 & k-means，シルエット係数\cite{rousseeuw1987} & 似た人のまとまりがあるか \\
4 & 構造 & 観測変数のパスモデル\cite{igolkina2020} & 理論上のつながりがデータと合うか \\
5 & 因果構造の探索 & PCアルゴリズム\cite{spirtes2000} & 矢印の向きが支持されるか \\
6 & 効果の推定 & Double Machine Learning\cite{chernozhukov2018} & 各行動が推し疲れ・幸福度をどれだけ動かすか \\
7 & 予測 & 線形回帰，10分割交差検証 & 見込みの幅 \\
8 & 群の比較 & ANCOVA，傾向スコア\cite{rosenbaum1983,hoshino2009} & 第5章の分析計画 \\
\bottomrule
\end{tabular}
\end{table}
